\documentclass[11pt]{article}
\usepackage[margin=1in]{geometry}
% \usepackage{booktabs}
\usepackage{array}
% \usepackage{longtable}
% \usepackage{caption}
% \usepackage{listings}
\usepackage[T1]{fontenc}
% \usepackage{inconsolata}
% \usepackage{hyperref}
% \hypersetup{hidelinks}
\setlength{\parskip}{0.6em}
\setlength{\parindent}{0pt}

\title{ECE6140 Project: Combinational Circuit Simulator}
\author{\vspace{-0.5em}}
\date{\vspace{-1.5em}}

\begin{document}
\maketitle

\section*{1. Data Structures (one page max)}
The simulator models a combinational circuit using the following core abstractions:

\textbf{Gate} (struct/class): represents a logic element. Each gate stores:
\begin{itemize}
  \item \texttt{gate\_name}: unique identifier (e.g., \texttt{and\_n10}).
  \item \texttt{gate\_type}: enum over \texttt{INPUT, OUTPUT, BUF, NOT, AND, OR, NAND, NOR, XOR, XNOR}.
  \item \texttt{input\_number}: arity derived from the netlist line.
  \item \texttt{gate\_input\_values}: vector of symbolic values \texttt{GateValue}.
  \item \texttt{gate\_output\_value}: a \texttt{GateValue} in \{0, 1, X, D, D'\}. (This enables five-valued logic; for this assignment we drive only 0/1.)
  \item Fan-in/out: \texttt{fan\_in\_gates} and \texttt{fan\_out\_gates} lists to traverse the graph.
\end{itemize}

\textbf{GateValue} (enum): \texttt{ZERO, ONE, UNKNOWN, D, D\_BAR} with overloaded logical operators (\texttt{\&}, \texttt{|}, \texttt{\textasciitilde}, \texttt{\^}) to evaluate multi-valued logic consistently across gates.

\textbf{Wire} (struct/class): names a net in the circuit and binds producers/consumers.
\begin{itemize}
  \item \texttt{wire\_name}: net label from the netlist.
  \item \texttt{input\_gate}: the single driver gate for this net.
  \item \texttt{output\_gates}: list of sinks receiving the net as input.
\end{itemize}

\textbf{Circuit} (container): owns all objects and provides APIs.
\begin{itemize}
  \item Maps: \texttt{wires: map<string, Wire>}, \texttt{gates: map<string, Gate>}.
  \item Port order: \texttt{input\_port\_names} and \texttt{output\_port\_names} (stable ordering for vector I/O).
  \item Build: \texttt{add\_a\_gate\_with\_wires(...)} creates gates, ensures wires exist, and records fan-in/out via wires.
  \item Connect: \texttt{connect\_gates()} walks wires to connect producers to consumers.
  \item Evaluate/reset: \texttt{evaluate()} performs a forward traversal; \texttt{reset\_gates()} clears state between vectors.
\end{itemize}

\textbf{CircuitWrapper}: thin helper to parse a netlist file (e.g., \texttt{s27.txt}) and populate a \texttt{Circuit}. It also provides \texttt{evaluate\_circuit\_with\_normal\_input(bitstring)} which:
\begin{enumerate}
  \item Resets all gate values to \texttt{UNKNOWN}.
  \item Assigns \texttt{ZERO/ONE} to each \texttt{INPUT} gate based on the input bitstring order.
  \item Calls \texttt{Circuit.evaluate()} then collects \texttt{OUTPUT} gate values in the declared output order.
\end{enumerate}

This design keeps the netlist graph explicit (via wires) and enables deterministic vector I/O ordering required for tabulated results.

\section*{2. Simulation Algorithm (pseudocode)}
The simulator uses event-driven forward propagation from inputs. Pseudocode below summarizes the core flow:

\begin{verbatim}
function evaluate_circuit_with_normal_input(input_bits):
    circuit.reset_gates()
    assert len(input_bits) == len(circuit.input_port_names)
    for i, input_name in enumerate(circuit.input_port_names):
        gate = circuit.gates[input_name]
        gate.set_output_value(ONE if input_bits[i] == '1' else ZERO)

    # Initialize evaluation queue with immediate fanout of inputs
    Q = deque()
    for input_name in circuit.input_port_names:
        Q.extend(circuit.gates[input_name].fan_out_gates)

    while Q not empty:
        g = Q.pop_left()
        prev = g.gate_output_value
        new  = g.forward()        # reads fan-in outputs, applies gate logic
        if prev == new:
            continue             # stabilized; skip re-enqueuing
        assert prev == UNKNOWN   # acyclic assumption for these benchmarks
        for h in g.fan_out_gates:
            Q.append(h)

    # Read outputs in the netlist-declared order
    out_bits = []
    for out_name in circuit.output_port_names:
        v = circuit.gates[out_name].gate_output_value
        out_bits.append('1' if v == ONE else '0' if v == ZERO else 'X')
    return ''.join(out_bits)
\end{verbatim}

\section*{3. Simulation Results (inputs and outputs)}
All vectors below were produced by running the provided implementation on the four benchmark circuits. Each table shows input and output bit-vectors in order.

\subsection*{s27}
\begin{center}
\small
\begin{tabular}{>{\ttfamily}l >{\ttfamily}l}
\hline
Input (7 bits) & Output (4 bits) \\
\hline
1110101 & 1001 \\
0001010 & 0100 \\
1010101 & 1001 \\
0110111 & 0001 \\
1010001 & 1001 \\
\hline
\end{tabular}
\end{center}

\subsection*{s298f\_2}
\begin{center}
\small
\begin{tabular}{>{\ttfamily}l >{\ttfamily}l}
\hline
Input (17 bits) & Output (20 bits) \\
\hline
10101010101010101 & 00000010101000111000 \\
01011110000000111 & 00000000011000001000 \\
11111000001111000 & 00000000001111010010 \\
11100001110001100 & 00000000100100100101 \\
01111011110000000 & 11111011110000101101 \\
\hline
\end{tabular}
\end{center}

\subsection*{s344f\_2}
\begin{center}
\small
\begin{tabular}{>{\ttfamily}l >{\ttfamily}l}
\hline
Input (24 bits) & Output (26 bits) \\
\hline
101010101010101011111111 & 10101010101010101010101101 \\
010111100000001110000000 & 00011110000000100001111100 \\
111110000011110001111111 & 00011100000111011000111010 \\
111000011100011000000000 & 00001101111001111111000010 \\
011110111100000001111111 & 10011101111000001001000100 \\
\hline
\end{tabular}
\end{center}

\subsection*{s349f\_2}
\begin{center}
\small
\begin{tabular}{>{\ttfamily}l >{\ttfamily}l}
\hline
Input (24 bits) & Output (26 bits) \\
\hline
101010101010101011111111 & 10101010101010101101010101 \\
010111100000001110000000 & 00011110000000101011110000 \\
111110000011110001111111 & 00011100000111010001111100 \\
111000011100011000000000 & 00001101111001110010001111 \\
011110111100000001111111 & 10011101111000001010000100 \\
\hline
\end{tabular}
\end{center}

\vspace{0.5em}
\textit{Reproducibility:} Run \texttt{python p1-1.py} in the provided environment; it builds each circuit from \texttt{./files/*.txt} and prints the vectors used to populate the tables above.

\end{document}


